\section{Kỹ năng mềm}

Trong quá trình thực hiện đồ án, bên cạnh những kiến thức kỹ thuật, các thành viên trong nhóm cũng đã có cơ hội rèn luyện và áp dụng hiệu quả nhiều kỹ năng mềm quan trọng. Những kỹ năng này đóng vai trò then chốt trong việc duy trì tiến độ, đảm bảo chất lượng công việc và tạo ra một môi trường làm việc nhóm tích cực.

\subsection{Kỹ năng giao tiếp và thảo luận}
Giao tiếp là yếu tố cốt lõi giúp nhóm kết nối và đồng bộ thông tin một cách xuyên suốt.
\begin{itemize}
    \item \textbf{Kênh liên lạc đa dạng:} Nhóm đã kết hợp sử dụng nhiều nền tảng như Zalo và Slack để trao đổi thông tin hằng ngày, giải đáp nhanh các thắc mắc và gửi các tài liệu tham khảo. 
    \item \textbf{Họp định kỳ:} Các buổi họp nhóm hàng tuần được tổ chức để cập nhật tiến độ, báo cáo các hạng mục đã hoàn thành và lên kế hoạch cho tuần tiếp theo. Qua đó, mỗi thành viên đều nắm rõ được định hướng chung của dự án.
    \item \textbf{Trình bày ý tưởng:} Khi đưa ra các giải pháp mới (ví dụ: áp dụng thuật toán 7-Bag Randomizer hay xử lý Wall Kicks), các thành viên đều trình bày rõ ràng ưu và nhược điểm để cả nhóm cùng đánh giá và đưa ra quyết định cuối cùng.
\end{itemize}

\subsection{Kỹ năng làm việc nhóm}
Một dự án phần mềm không thể hoàn thiện nếu thiếu đi sự phối hợp ăn ý giữa các thành viên. Nhóm đã áp dụng kỹ năng làm việc nhóm một cách triệt để:
\begin{itemize}
    \item \textbf{Phân chia công việc hợp lý:} Công việc được chia nhỏ và giao phó dựa trên điểm mạnh của từng người. Ví dụ, những thành viên có nền tảng tư duy thuật toán vững sẽ đảm nhiệm logic cốt lõi của game (xoay khối, xóa dòng, hệ thống điểm), trong khi những bạn có thế mạnh về tổng hợp thông tin sẽ tập trung vào việc soạn thảo tài liệu kỹ thuật, báo cáo quá trình và slide thuyết trình.
    \item \textbf{Tương trợ lẫn nhau:} Các thành viên không làm việc một cách biệt lập. Khi một cá nhân gặp "blocker" (trở ngại) trong quá trình code hoặc gặp lỗi logic khó tìm (bug), các thành viên khác sẵn sàng tham gia debug chéo hoặc review code để cùng giải quyết vấn đề.
    \item \textbf{Tinh thần trách nhiệm:} Tất cả các thành viên đều thể hiện tính kỷ luật cao, cam kết hoàn thành đúng deadline các thẻ công việc (Card) đã được giao trên Trello, từ trạng thái "To Do" sang "Done".
\end{itemize}

\subsection{Kỹ năng giải quyết xung đột và khó khăn}
Trong một dự án thực tế, những ý kiến trái chiều hay các khó khăn kỹ thuật là điều khó tránh khỏi. Nhóm đã giải quyết các vấn đề này thông qua các nguyên tắc sau:
\begin{itemize}
    \item \textbf{Xử lý mâu thuẫn kỹ thuật:} Khi có sự bất đồng về cách thiết kế kiến trúc hoặc cách triển khai một tính năng (ví dụ: cách kế thừa class, cách tổ chức thư mục code), nhóm luôn lấy tiêu chí "tối ưu và dễ bảo trì" làm trọng tâm. Các quyết định được đưa ra dựa trên việc thảo luận minh bạch và biểu quyết đa số.
    \item \textbf{Xử lý xung đột mã nguồn (Merge Conflict):} Khi làm việc chung trên một Repository, việc bị conflict code là rất thường xuyên xảy ra. Thay vì tự ý ghi đè code của người khác, các thành viên sẽ chủ động liên hệ với người viết đoạn code đó để cùng nhau hợp nhất (merge) một cách an toàn nhất, đảm bảo tính năng của cả hai đều hoạt động tốt.
\end{itemize}
